\chapter{Wartości własne, diagonalizacja i twierdzenie spektralne}
\label{ch:lin14-spektralne}
\index{wartość własna@Wartość własna}
\index{wektor własny@Wektor własny}
\index{diagonalizacja@Diagonalizacja}
\index{operator normalny@Operator normalny}
\index{operator samosprzężony@Operator samosprzężony}
\index{twierdzenie spektralne@Twierdzenie spektralne}

\noindent\textbf{Cel rozdziału.}
Rozdziały~\ref{ch:lin12-euklidesowe-unitarne}
i~\ref{ch:lin13-ortogonalnosc-rzuty}
wprowadziły iloczyny skalarne, operator sprzężony
i rzuty ortogonalne. Teraz badamy operator liniowy
za pomocą jego wartości własnych. Najpierw omówimy
diagonalizację bez założenia istnienia iloczynu skalarnego.
Następnie udowodnimy twierdzenie Schura nad
\(\mathbb C\) i wyprowadzimy twierdzenie spektralne
dla operatorów normalnych, hermitowskich i
rzeczywistych samosprzężonych. Sformułujemy też
rozkład na rzuty spektralne i jego konsekwencje.

Przez \(K\) oznaczamy dowolne ciało, dopóki nie
wprowadzimy jawnie założenia \(K=\mathbb R\)
lub \(K=\mathbb C\).
W każdej części tego rozdziału \(V\) jest przestrzenią
skończenie wymiarową, a \(T\colon V\to V\)
endomorfizmem liniowym. Wymiar \(n=\dim_K V\)
może być równy zeru.
Przy iloczynie skalarnym stosujemy konwencję
liniowości w drugim argumencie, jak w
Rozdziale~\ref{ch:lin12-euklidesowe-unitarne}.

\section{Wartości własne i wielomian charakterystyczny}
\label{sec:lin14-wartosci}

\begin{definicja}[Wartość własna i przestrzeń własna]
\label{def:lin14-wlasne}
Niech \(T\in\operatorname{End}_K(V)\).
Skalar \(\lambda\in K\) nazywamy
\emph{wartością własną} \(T\),
jeśli istnieje \(v\in V\setminus\{0\}\) taki, że
\[
 Tv=\lambda v.
\]
Każdy taki \(v\) jest \emph{wektorem własnym}
odpowiadającym \(\lambda\).
\emph{Przestrzenią własną} dla \(\lambda\) jest
\(E_\lambda(T):=\ker(T-\lambda I_V)\).
Zbiór wartości własnych w \(K\) oznaczamy
\(\sigma_K(T)\).
Wartość własna musi mieć
niezerową przestrzeń własną.
\end{definicja}

\begin{definicja}[Wielomian charakterystyczny]
\label{def:lin14-charpoly}
Niech \(\dim_K V=n\) i \(B\) będzie
uporządkowaną bazą \(V\).
\emph{Wielomianem charakterystycznym}
\(T\) nazywamy
\[
 \chi_T(t):=\det(tI_n-[T]_B)\in K[t].
\]
Dla \(n=0\) przyjmujemy \(\chi_T=1\).
Jeżeli \(\lambda\in K\) jest pierwiastkiem
\(\chi_T\), liczbę jego wystąpień
z krotnością nazywamy
\emph{krotnością algebraiczną}
\(a_T(\lambda)\).
\emph{Krotnością geometryczną}
nazywamy
\(g_T(\lambda):=\dim_K E_\lambda(T)\).
\end{definicja}

\begin{stwierdzenie}[Niezależność bazy i kryterium wartości własnej]
\label{prop:lin14-charpoly}
Niech \(V\) będzie skończenie wymiarowa nad
ciałem \(K\), \(T\in\operatorname{End}_K(V)\),
a \(\lambda\in K\).
Wielomian \(\chi_T\) nie zależy
od wybranej bazy, ma stopień
\(n=\dim V\) i współczynnik wiodący \(1_K\).
Następujące warunki są równoważne:
\begin{enumerate}[label=(\roman*)]
\item \(\lambda\) jest wartością własną \(T\);
\item \(T-\lambda I_V\) nie jest iniektywne;
\item \(T-\lambda I_V\) nie jest odwracalne;
\item \(\chi_T(\lambda)=0_K\).
\end{enumerate}
\end{stwierdzenie}
\begin{proof}
Niech \(A=[T]_B\).
Jeżeli \(B'\) jest inną bazą, a
\(P\) macierzą przejścia od \(B'\) do \(B\),
to
\([T]_{B'}=P^{-1}AP\).
Stąd
\[
 \det(tI-P^{-1}AP)
 =\det(P^{-1}(tI-A)P)=\det(tI-A).
\]
Współczynnik najwyższej potęgi \(t^n\)
w wyznaczniku \(tI-A\) powstaje
z wybrania wszystkich składników \(t\)
na przekątnej i ma wartość \(1_K\);
pozostałe składniki mają niższy stopień.
Przy \(n=0\) jest to wielomian stały \(1_K\).

Warunki (i) i (ii) są równoważne
na mocy definicji wektora własnego.
Ponieważ \(V\) ma skończony wymiar,
endomorfizm jest iniektywny dokładnie
wtedy, gdy jest odwracalny;
to daje (ii)\(\Leftrightarrow\)(iii).
Z kryterium odwracalności macierzy
przez wyznacznik otrzymujemy
\[
 T-\lambda I_V\text{ nieodwracalne}
 \ \Longleftrightarrow\
 \det(A-\lambda I_n)=0_K
 \ \Longleftrightarrow\
 \chi_T(\lambda)=0_K.
\]
Przy \(n=0\) wszystkie cztery warunki
są fałszywe, zgodnie z \(\chi_T=1\).
\end{proof}

\begin{stwierdzenie}[Granice krotności geometrycznej]
\label{prop:lin14-krotnosci}
Niech \(T\in\operatorname{End}_K(V)\)
i \(\lambda\in K\) będzie jego wartością własną.
Wtedy
\[
 1\le g_T(\lambda)\le a_T(\lambda).
\]
\end{stwierdzenie}
\begin{proof}
Dolna granica wynika z istnienia
niezerowego wektora własnego.
Niech \(g=g_T(\lambda)\).
Wybierzmy bazę
\((e_1,\ldots,e_g)\) przestrzeni \(E_\lambda(T)\)
i uzupełnijmy ją do bazy \(V\).
Ponieważ \(Te_j=\lambda e_j\)
dla \(1\le j\le g\), macierz \(T\)
w takiej bazie ma postać blokową
\[
 [T]=\begin{pmatrix}
 \lambda I_g&C\\0&D
 \end{pmatrix}.
\]
Wyznacznik macierzy blokowo trójkątnej daje
\[
 \chi_T(t)=(t-\lambda)^g\det(tI_{n-g}-D).
\]
Wielomian jest więc podzielny
przez \((t-\lambda)^g\),
czyli \(a_T(\lambda)\ge g\).
Gdy \(g=n\), drugi wyznacznik
jest wyznacznikiem macierzy pustej
równym \(1_K\).
\end{proof}

\begin{przyklad}[Ta sama wartość własna, różna krotność geometryczna]
\label{prz:lin14-jordan2}
Nad dowolnym ciałem \(K\) rozważmy
\[
 A=\begin{pmatrix}1&1\\0&1\end{pmatrix},
 \qquad
 D=\begin{pmatrix}1&0\\0&1\end{pmatrix}.
\]
Obie macierze mają
\(\chi(t)=(t-1)^2\), zatem
wartość własna \(1\) ma krotność
algebraiczną \(2\).
Jednak
\[
 \ker(A-I)=\operatorname{span}\{(1,0)\},
 \qquad
 \ker(D-I)=K^2.
\]
Stąd \(g_A(1)=1\) i \(g_D(1)=2\).
Pierwsza macierz nie może być
diagonalizowana, druga już jest diagonalna.
Sam wielomian charakterystyczny nie
rozstrzyga diagonalizowalności.
\end{przyklad}

\section{Diagonalizacja i niezależność przestrzeni własnych}
\label{sec:lin14-diagonalizacja}

\begin{twierdzenie}[Niezależność wektorów o różnych wartościach własnych]
\label{thm:lin14-eigen-independence}
Niech \(T\in\operatorname{End}_K(V)\),
a \(\lambda_1,\ldots,\lambda_r\in K\)
będą parami różnymi wartościami własnymi.
Jeśli \(v_j\in E_{\lambda_j}(T)\setminus\{0\}\),
to \(v_1,\ldots,v_r\) są liniowo niezależne.
Co więcej,
\[
 E_{\lambda_1}(T)+\cdots+
 E_{\lambda_r}(T)
 =
 E_{\lambda_1}(T)\oplus\cdots\oplus
 E_{\lambda_r}(T).
\]
\end{twierdzenie}
\begin{proof}
Załóżmy, że
\(v_j\in E_{\lambda_j}(T)\)
oraz \(\sum_{j=1}^rv_j=0\).
Dla ustalonego \(k\in\{1,\ldots,r\}\)
zastosujmy do tej równości
odwzorowanie liniowe
\[
 R_k=\prod_{j\ne k}(T-\lambda_jI_V).
\]
Czynniki przemieniają się, ponieważ
wszystkie są wielomianami w \(T\).
Jeśli \(j\ne k\), w iloczynie znajduje się
\(T-\lambda_jI_V\), który zeruje \(v_j\);
zatem \(R_kv_j=0\).
Natomiast
\[
 R_kv_k=
 \prod_{j\ne k}(\lambda_k-\lambda_j)v_k.
\]
Każdy czynnik skalarny jest niezerowy;
w ciele ich iloczyn również jest niezerowy.
Wynika stąd \(v_k=0\).
Ponieważ \(k\) był dowolny,
wszystkie \(v_j\) są zerowe.
To dowodzi bezpośredniości sumy
przestrzeni własnych i, po wybraniu
po jednym niezerowym wektorze z każdej
z nich, pierwszej tezy.
\end{proof}

\begin{definicja}[Operator diagonalizowalny]
\label{def:lin14-diagonalizable}
Endomorfizm \(T\in\operatorname{End}_K(V)\)
nazywamy \emph{diagonalizowalnym nad \(K\)},
jeśli istnieje baza \(V\)
złożona z wektorów własnych \(T\).
Równoważnie jego macierz
w pewnej bazie jest diagonalna.
Nie zakładamy istnienia
iloczynu skalarnego.
\end{definicja}

\begin{twierdzenie}[Kryterium diagonalizowalności]
\label{thm:lin14-diag-criterion}
Niech \(T\) będzie endomorfizmem
przestrzeni \(V\) nad \(K\),
o wymiarze \(n<\infty\).
Następujące warunki są równoważne:
\begin{enumerate}[label=(\roman*)]
\item \(T\) jest diagonalizowalny nad \(K\);
\item \(V=\bigoplus_{\lambda\in\sigma_K(T)}
                    E_\lambda(T)\);
\item \(\sum_{\lambda\in\sigma_K(T)}
             g_T(\lambda)=n\);
\item wielomian charakterystyczny
      rozkłada się nad \(K\) na czynniki liniowe
      i dla każdej wartości własnej
      \(g_T(\lambda)=a_T(\lambda)\).
\end{enumerate}
W szczególności jeśli \(\chi_T\)
ma \(n\) parami różnych pierwiastków w \(K\),
to \(T\) jest diagonalizowalny.
Przy \(n=0\) wszystkie warunki zachodzą
z odpowiednimi pustymi sumami.
\end{twierdzenie}
\begin{proof}
Jeśli istnieje baza złożona
z wektorów własnych, pogrupowanie
jej elementów według wartości własnych
daje bazę sumy odpowiednich
przestrzeni własnych.
Suma jest bezpośrednia na mocy
Twierdzenia~\ref{thm:lin14-eigen-independence},
co dowodzi (i)\(\Rightarrow\)(ii).
Jeśli zachodzi (ii), wybór bazy
w każdej z przestrzeni własnych
i połączenie tych baz daje
bazę \(V\) złożoną z wektorów własnych;
zatem (ii)\(\Rightarrow\)(i).
Ponieważ suma jest bezpośrednia
niezależnie od tego, czy wypełnia \(V\),
jej wymiar jest sumą
wymiarów składników.
Wobec tego (ii)\(\Leftrightarrow\)(iii).

Dla (i)\(\Rightarrow\)(iv)
w bazie własnej mamy
macierz diagonalną
\(\operatorname{diag}(\lambda_1,\ldots,\lambda_n)\);
stąd \(\chi_T(t)=\prod_{j=1}^n(t-\lambda_j)\).
Dla ustalonej wartości własnej
liczba wystąpień na przekątnej
jest wymiarem jej przestrzeni własnej,
bo wektor własny o wartości \(\lambda\)
ma współrzędne zerowe
przy wszystkich innych wartościach.
Zatem \(g_T(\lambda)=a_T(\lambda)\).

Jeśli zachodzi (iv), to z
rozszczepienia wielomianu
\(\chi_T\) i jego stopnia \(n\)
otrzymujemy
\[
 n=\sum_{\lambda\in\sigma_K(T)}
                         a_T(\lambda)
  =\sum_{\lambda\in\sigma_K(T)}
                         g_T(\lambda).
\]
Jest to (iii).
Jeśli pierwiastków jest \(n\)
i wszystkie są różne,
każda z ich przestrzeni własnych
ma wymiar przynajmniej \(1\).
Suma bezpośrednia tych przestrzeni
ma więc wymiar co najmniej \(n\);
nie może mieć więcej niż \(n\),
zatem wypełnia \(V\).
\end{proof}

\begin{przyklad}[Diagonalizacja bez iloczynu skalarnego]
\label{prz:lin14-diagonalization2}
Niech
\[
 A=\begin{pmatrix}2&1\\0&3\end{pmatrix}
 \quad\text{nad }\mathbb R.
\]
Wielomian charakterystyczny jest
równy \((t-2)(t-3)\).
Wektor \(v_1=(1,0)^{\mathsf T}\)
ma wartość własną \(2\), a
\(v_2=(1,1)^{\mathsf T}\)
ma wartość własną \(3\).
Wybierając
\[
 P=\begin{pmatrix}1&1\\0&1\end{pmatrix},
 \qquad
 D=\begin{pmatrix}2&0\\0&3\end{pmatrix},
\]
sprawdzamy \(AP=PD\),
czyli \(P^{-1}AP=D\).
Macierz \(P\) nie musi być ortogonalna:
diagonalizowalność ogólnego operatora
jest pojęciem algebraicznym,
a nie własnością zależną od iloczynu skalarnego.
\end{przyklad}
